%=== CHAPTER SIX ===
%% Rewritten 2026-10-30 against the 14-state mass-only interface; 2026-11-05:
%% objective 2 distinguishes the full-state (campaign) and reduced (deployed) estimator.
%% 2026-11-03: findings reorganised by research objective (answer / evidence /
%% boundary / not established); future work reduced to four ranked items.
%% 2026-11-05: release-campaign results (312 flights) added to objectives 3-5.
%% 2026-11-07: findings led by the 180-flight deployed-system validation
%% (Section sec:eval:deployed); early confirmation under an under-reading thrust
%% map with delayed detachment added as a limit; future work re-ranked.

\chapter{Conclusion and Future Work}
\label{ch:conclusion}
\begin{spacing}{1.5}
\setlength{\parskip}{0.22in}

\section{Summary}

This dissertation treated payload grasp and release on a quadrotor as one
estimation-and-control problem. Chapter~\ref{ch:model} reduced a centred rigid
payload to a single online parameter, the composite mass, from which the roll
and pitch inertia increment follows, and assigned the identification of that
mass to the translational channel. Chapter~\ref{ch:framework} built a
moving-horizon estimator driven by a wrench reconstructed from rotor speeds, a
predictive controller updated continuously from one health-gated estimate, and a
release lifecycle in which acceptance of a release is gated on evidence --- in
the deployed form an independent force balance, in the historical cohort the
estimator's mass evidence. Chapter~\ref{ch:evaluation}
evaluated the deployed system --- the reduced estimator with the force-balance
rule --- in a validation campaign of $180$ PX4/Gazebo software-in-the-loop
flights, and documented how the design was reached with historical batches flown
with a full-state estimator: a frozen cohort of $120$ flight attempts, a smaller
frozen ablation, a release campaign of $312$ flights and a paired estimator
comparison of $96$. The findings below answer the five research objectives
against that evidence, and say which batch each statement rests on.

\section{Principal Findings}

\paragraph{Objective 1 --- a single-parameter payload structure.}
\begin{itemize}
\item \textbf{Answer.} For a rigid payload on the thrust axis, the roll and pitch
inertia increment is $\mu r_z^2$, a function of the composite mass and a nominal
arm, so the composite mass is the only payload quantity that needs online
identification; the inertia follows from it and vanishes with it.
\item \textbf{Evidence.} The derivation of Section~\ref{sec:model:payload}; the
structure is implemented in both optimisers and used in every reported flight.
\item \textbf{Boundary.} A centred grasp within the alignment tolerance, a
point-mass payload, and a nominal arm whose realised value varies by up to
$-14$ to $+23$ percent in inertia (Section~\ref{sec:model:approx}).
\item \textbf{Not established.} Eccentric grasps; and whether estimating the
inertia directly would perform better or worse, which was not compared.
\end{itemize}

\paragraph{Objective 2 --- identifying the mass.}
\begin{itemize}
\item \textbf{Answer.} The estimator identifies the composite mass accurately in
loaded hover, but under sustained manoeuvre the estimate is biased low and
periodically collapses toward the bare airframe. Both statements describe the
full-state estimator with which the campaigns were flown. The deployed reduced
estimator, which does not integrate the rotational dynamics, removes most of
that bias and the collapse, in replay and in closed loop, at about a tenth of
the computation. The improvement belongs to its model structure --- reconstructed
thrust and measured attitude as known inputs, no rotational channel --- rather
than to the moving horizon: replayed on the same flights, an extended Kalman
filter with that structure was as accurate.
\item \textbf{Evidence.} In the deployed-system validation the reduced
estimator's loaded mass error was \SI{0.013}{\kg} at the median over $72$
nominal releases (\SI{2.7}{\gram} in loaded hover, a figure-eight bias of
\SI{-0.66}{\percent} at \SI{4}{\meter\per\second}), its estimate never
collapsed in any flight, it never re-initialised, and it solved in
\SI{0.31}{\milli\second} at the median (Section~\ref{sec:eval:deployed}). In
the historical frozen cohort the per-flight loaded-hover error
lay between \SI{-0.44}{\percent} and \SI{-0.10}{\percent}; the loaded
figure-eight bias had medians of \SI{-1.0}{\percent} at
\SI{2}{\meter\per\second} and \SI{-4.1}{\percent} and \SI{-3.3}{\percent} at
\SI{4}{\meter\per\second}, with dips to the mass lower bound and a median
collapse fraction of $0.38$ for the lighter payload
(Section~\ref{sec:eval:estimation}). Replayed on the 312 campaign flights, the reduced
estimator lowered the median loaded mass error from \SI{0.036}{\kg} to
\SI{0.013}{\kg}, smaller in every one of the 264 nominal flights, and returned to
the empty value after every release, where the full-state estimate stayed about
\SI{2}{\percent} low in 57 flights at \SI{4}{\meter\per\second}. In closed loop,
over $48$ pairs flown with the same controller and release rule, the reduced
estimator lowered the loaded mass error from \SI{0.035}{\kg} to \SI{0.013}{\kg}
(better in $46$ pairs) and the post-release error from \SI{0.030}{\kg} to
\SI{0.010}{\kg} (better in all $48$); at \SI{4}{\meter\per\second} the full-state
estimate spent a median \SI{22}{\percent} of the loaded figure-eight below half
the payload and the reduced estimate none, in any flight. The median solve time
fell from \SI{2.8}{\milli\second} to \SI{0.3}{\milli\second}. Tracking and the
lifecycle did not change. In replay on the $48$ reduced-estimator flights, the
extended Kalman filter came within \SI{0.2}{\gram} of the reduced estimator's loaded
mass error, converging \SI{0.5}{\second} later after lift-off and
\SI{0.4}{\second} sooner after detachment; the window force balance was
\SIrange{0.5}{1.6}{\gram} less accurate. Penalising the newest measurement, which
the flown estimator leaves without cost, shortened its convergence by
\SIrange{0.2}{0.3}{\second} at unchanged accuracy, and replacing the body-origin
velocity by the composite-centre velocity changed nothing
(Section~\ref{sec:eval:reduced}).
\item \textbf{Boundary.} Simulated rotor speeds are exact and the plant uses the
thrust law the estimator inverts; the shortened arm; payloads of
\SIrange{0.15}{0.30}{\kg}.
\item \textbf{Not established.} The full cause of the collapse; the reduced
estimator under realistic attitude noise, which the simulator's odometry does not
exercise; the effect of a thrust-coefficient error on the actuation side, and of
an error smaller than \SI{10}{\percent} (Section~\ref{sec:eval:thrustcal});
whether the moving horizon is needed at all, since the filter baseline was
compared only in open-loop replay with nominal thrust; and the advantage of the
window seed, whose comparison was not run.
\end{itemize}

\paragraph{Objective 3 --- a continuously updated controller.}
\begin{itemize}
\item \textbf{Answer.} The controller completed every eligible flight without
divergence while consuming the estimate continuously. The residual specific force
improves tracking but does not improve the mass estimate or change the
lifecycle. The held inertia floor protects the inner loop through the collapse,
but for that reason the cohort cannot show the effect of estimate-driven
inner-loop scheduling on its own.
\item \textbf{Evidence.} None of the $180$ flights of the deployed-system
validation diverged or logged a controller solve failure, including the jammed,
lost-load and thrust-map flights (Section~\ref{sec:eval:deployed}). No eligible
cohort flight diverged (Table~\ref{tab:attempts}). In the $48$-flight ablation the residual reduced
loaded figure-eight RMSE by \SIrange{22}{31}{\percent} and removed a mean
vertical sag of up to \SI{44}{\milli\meter}, with no difference in mass bias,
confirmation delay or carry-phase confirmations
(Section~\ref{sec:eval:l1ablation}). The inner-loop scale stayed at the floor
value of $2.34$ through every loaded figure-eight
(Section~\ref{sec:eval:opslimit}). In the release campaign, with the lifecycle
and inner loop shared, the MHE controller had a smaller loaded RMSE than a
fixed-parameter controller in all $24$ pairs, by \SIrange{0.04}{0.14}{\meter}
at the median, and no consistent difference in carry from an $\mathcal{L}_1$
controller that uses the residual alone (Section~\ref{sec:eval:camp:ctrl}).
\item \textbf{Boundary.} The inertia ratio at the shortened arm stays below the
scaling cap; with the earlier, longer arm it would not.
\item \textbf{Not established.} A tracking benefit of estimating the mass over
the lumped residual alone --- the only consistent difference was a smaller
post-release peak at \SI{2}{\meter\per\second}, of \SIrange{6}{17}{\milli\meter};
any benefit over incremental dynamic inversion or full $\mathcal{L}_1$-adaptive
NMPC, which were not run; and the contribution of the continuous inner-loop
scaling, which the held floor masked during carry and which all three
controllers of the comparison shared.
\end{itemize}

\paragraph{Objective 4 --- separating a release command from a physical
detachment.}
\begin{itemize}
\item \textbf{Answer.} Gating release acceptance on evidence separates
``the command was issued'' from ``the payload has left''. On the mass evidence,
only consulting it after the command --- \emph{command-armed,
estimator-confirmed detachment} --- is sound; consulting it throughout carry is
not, because the mass evidence fails in manoeuvre. With independent
force-balance evidence, which does not pass through the mass estimate, the rule
can in addition be enabled by a watchdog armed on airborne load: it then also
detects a load lost without a command, without false acceptances in carry. That
rule's weakness is the thrust map, not the manoeuvre: when the map under-reads by
more than the payload's share of the composite mass, a loaded vehicle reads as
empty, and the rule then accepts a delayed release before the payload has left.
\item \textbf{Evidence.} In the deployed-system validation the force-balance
rule confirmed $72$ of $72$ nominal releases, none early, none during carry,
\SI{2.17}{\second} after physical detachment at the median; it detected $24$ of
$24$ uncommanded losses, \SI{2.15}{\second} after detachment, and accepted none
of $24$ jammed releases, each of which ended in a braked hover. With a
\SI{10}{\percent} over-reading thrust map it missed all $24$ releases and hovered
safely. With a \SI{10}{\percent} under-reading map the \SI{0.15}{\kg} payload,
\SI{6.8}{\percent} of the composite mass, read as absent: the loss watchdog never
armed, and with detachment held \SI{4}{\second} the release was accepted in $12$
of $12$ flights, \SI{3.1}{\second} before the payload left
(Section~\ref{sec:eval:deployed}). In the historical cohort every eligible release was accepted,
\SIrange{2.3}{3.2}{\second} after physical detachment. Under delayed detachment
the command-triggered arm accepted early in $24$ of $24$ flights, about
\SI{4}{\second} early, and the evidence-gated arms in none of $47$. The
always-on rule accepted $114$ releases that had not happened, in $20$ of $32$
flights at \SI{4}{\meter\per\second}; the command-armed rule accepted none in its
$23$ flights and confirmed real detachments no later
(Sections~\ref{sec:eval:completion}--\ref{sec:eval:unsignaled}). In the
historical release campaign the force-balance rule accepted no release during carry in any
of $240$ flights ($228$ with the watchdog armed), none early in $36$
delayed-detachment flights, against $36$ of $36$ for command-triggered
acceptance, and confirmed real detachments \SIrange{1.9}{2.5}{\second} after
they occurred. It detected $24$ of $24$ uncommanded losses
\SIrange{1.8}{2.5}{\second} after detachment, resetting the payload geometry
and inner-loop scale each time, and accepted none of $24$ releases with a jammed
gripper, all of which ended in a braked hover (Section~\ref{sec:eval:campaign}).
\item \textbf{Boundary.} For the command-armed rule, a release command must
exist. For both rules, the detachment delay tested was \SI{4}{\second}. The
force-balance rule assumes a thrust map accurate to better than
\SI{10}{\percent} and, more precisely, accurate to better than the payload's
share of the composite mass: overestimated by \SI{10}{\percent} it missed every
release, safely; underestimated by as much it read the lighter payload as absent,
so that the loss watchdog never armed and a delayed release was accepted early
(Sections~\ref{sec:eval:deployed} and~\ref{sec:eval:camp:thrust}). For the
deployed system this is a calibration requirement on $k_f$, not a tuning
choice.
\item \textbf{Not established.} A tracking or survival advantage of
evidence-gated acceptance --- none was detected, and none was expected, since the
model's mass follows the estimate in every arm; what the delayed-detachment
experiment shows is lifecycle correctness. Also not established: uncommanded
loss under the mass-evidence rules, which was not run; the force-balance rule
under delayed detachment with an over-reading map or with the heavier payload,
and under a thrust-map error smaller than \SI{10}{\percent}; detection of
partial loss; and its performance relative to an optimised change detector.
\end{itemize}

\paragraph{Objective 5 --- evaluation and its scope.}
\begin{itemize}
\item \textbf{Answer.} The deployed system was validated in a campaign of its
own, from a tagged commit, with randomised rounds, complete reporting and
acceptance criteria inherited from an earlier pre-specification; the historical
configuration was evaluated with randomised blocks, complete reporting of every
attempt, targeted ablations and an explicit account of its failure modes.
\item \textbf{Evidence.} The validation: $180$ flights, all eligible, none
rerun, the four inherited acceptance criteria met; the criteria were fixed during
the run, after $11$ flights, not before it (Section~\ref{sec:eval:dep:accept}).
The cohort: $120$ attempts, $118$ eligible, two exclusions both
explained (Table~\ref{tab:attempts}); the historical campaign: $312$ attempts, all
eligible, with no replacement flight needed (Section~\ref{sec:eval:camp:design});
the evidence sources and their standing are listed in Table~\ref{tab:evidence}.
\item \textbf{Boundary.} All system-level evidence is software-in-the-loop.
\item \textbf{Not established.} Hardware validity; a comparison with
incremental dynamic inversion or with a change detector under a matched
false-alarm constraint; and the designed small thrust-map, partial-loss and
window-seed experiments, which were not run, and the parts of the thrust-map and
delayed-detachment crossing that were not flown (Section~\ref{sec:eval:notrun}).
\end{itemize}

\section{Limitations}

The limitations that bound the findings are collected here once. All
system-level evidence is software-in-the-loop, with exact rotor speeds and a plant
that generates thrust by the same law the estimator inverts. The attachment is
rigid and the grasp centred; eccentric grasps, gripper compliance, pendulum
modes and flexible loads are not modelled, and a sustained vertical gust is
intrinsically confounded with a payload change on a collective-thrust residual.
Every flight uses one gripper geometry, whose nominal arm differs from the
realised one by an unmeasured amount. The attach-time inertia floor is armed
whether or not the grasp succeeds. Under the mass-evidence rules it is ended
only by the release command, so a load lost without a command leaves it in
place; under the force-balance rule a detected loss ends it, but detection
depends on the thrust map. The release command is not gated on estimator
health. The evaluated configuration includes settings that are not
the defaults of the campaign code base's experiment entry points, and the results
are claims about that configuration; in the mainline code base those settings are
the defaults, and the deployed system differs from the evaluated configuration
only in its release rule and its estimator (Section~\ref{sec:eval:config}). The
deployed system was validated in $180$ flights, which cover nominal and delayed
release, uncommanded loss, a jam and a \SI{\pm10}{\percent} estimation-side
thrust-map error; its acceptance criteria were fixed during that run, not before
it. Its decisive limit is the thrust map: under a \SI{10}{\percent}
under-reading map it cannot tell the lighter payload from none, so the loss
watchdog does not arm and a delayed release is accepted about \SI{3}{\second}
before the payload leaves (Section~\ref{sec:eval:deployed}).

Three further limits concern what the evidence can show. The force-balance
detector was validated in a plant whose thrust follows exactly the static law it
inverts, with near-ideal rotor speeds and an inertial unit at the body origin;
rotor-speed noise and latency, accelerometer bias and vibration, the lever arm
between the inertial unit and the composite centre of mass, induced inflow and
an actuation-side thrust error were not present, and the thrust-map error was
applied to the estimation side only. Repeated flights of one cell differ through
unseeded sensor noise and scheduling, not through randomised physical
conditions, so the statistical tests describe repeatability in simulation, not
inference to a fielded vehicle (Section~\ref{sec:eval:stats}). And the release
rules were compared with each other, never with an optimised change detector at
a matched false-alarm rate; that is the largest gap for the release question.
The controller's torque bounds are a surrogate, not the realisable input set,
and the torque is a virtual input that the autopilot never receives
(Section~\ref{sec:frame:inputs}).

\section{Future Work}

The following directions are ranked by how much they would change what can be
claimed.

\begin{enumerate}
\item \textbf{Release evidence that knows its own calibration.} The
validation showed that the deployed rule accepts a delayed release early when
the thrust map under-reads by more than the payload's share of the composite
mass. A rule that combines $m_{\mathrm{acc}}$ with an explicit bound on the
thrust-map uncertainty --- refusing to treat a vehicle as empty, or reporting
that it cannot tell, when the expected payload is smaller than that bound, and
setting the thresholds of \eqref{eq:gammaacc} from the calibrated uncertainty
rather than as absolute masses --- should be designed and tested against the
crossing of thrust-map error and delayed detachment, of which only one corner
was flown. Evidence insensitive to a constant coefficient error, such as the
relative step in $m_{\mathrm{acc}}$ against the value recorded after lift, or
an online calibration of $k_f$ in empty hover, belongs to the same test. None of
these has been implemented or validated.

\item \textbf{The estimator the deployed system needs.} The extended Kalman
filter that matched the reduced estimator in replay should be flown in closed
loop against it, under nominal and erroneous thrust maps, to settle whether a
moving horizon is needed; the same batch can add the terminal measurement to the
estimator and replace the controller's torque box by the four per-rotor force
limits.

\item \textbf{Manoeuvre-robust mass estimation.} The reduced estimator removed
the collapse that dominated the historical configuration, but a manoeuvre bias of
about \SI{-0.7}{\percent} remains at \SI{4}{\meter\per\second}, and the
reference-point correction does not explain it. Calibrating the
airframe's rotational model --- inertia, rotor-drag moment and torque bias ---
would reduce the error that reaches the mass through the attitude, and might
eventually let the rotational channel contribute information instead of error.
The designed but unrun partial-loss and window-seed experiments belong here: both
test the estimator rather than the lifecycle.

\item \textbf{Health-gated release.} The release command should not be
issued while the estimator is unhealthy, which the deployed system does not
enforce.

\item \textbf{Change-detection and rotational baselines.} A comparison against
CUSUM or a related procedure on the same signals, under a matched false-alarm
constraint, would show whether persistence on a filtered score is near the
achievable detection performance. Since the mass estimate gave no consistent
tracking benefit over the lumped residual in carry, its value should next be
tested where the residual cannot substitute for it --- in the rotational
channel, with the inner-loop scaling no longer shared between the arms.

\item \textbf{Thrust-map robustness and hardware transfer.} The one-for-one
sensitivity to the thrust coefficient was measured at $\pm\SI{10}{\percent}$,
for the release detector and the estimator together. It should next be measured
at the designed $\pm\SI{5}{\percent}$, with the error also applied to the
actuation side, and then against motor-stand calibration of $k_f$ and $k_m$. Transfer should proceed through
tethered hover with the estimator running open loop, centred pickup, and only
then dynamic flight, with electronic-speed-controller telemetry replacing
simulated rotor speeds and a supervisor gating every stage.
\end{enumerate}

\section{Closing Remark}

The result that generalises beyond this vehicle concerns the relation between an
estimate and a decision. A continuously estimated mass can drive the model of a
predictive controller without any event, and doing so kept every flight of the
frozen cohort and of the deployed-system validation stable. Deciding that a payload has gone is a different matter: the
same estimate that is accurate enough to adapt a model in hover is not reliable
enough to declare a release in manoeuvre. In this work that decision was safe
either when the estimate was asked a narrower question --- whether a release that
has already been commanded has physically happened --- or when it was not asked
at all, and the release question was put to a measurement that does not depend
on it. The second form also detects a load lost without a command; it moves the
weak point from the manoeuvre to the calibration of the thrust map, and the
validation located that weak point exactly: an under-reading map larger than the
payload's share of the vehicle's mass makes a loaded vehicle look empty.

\end{spacing}
\newpage
